\chapter{Threat Model and Adversarial Analysis}
\label{ch:threat}

\section{Scope of the Analysis}
\label{sec:threat-scope}

A verified toolchain is only as trustworthy as the assumptions it makes
explicit. This chapter states the adversary model for MATHLIB5, the
guarantees each boundary provides against that adversary, and the residuals
the system deliberately leaves out of scope.

\section{Adversary Capabilities}
\label{sec:threat-adv}

We assume an adversary that can:
\begin{itemize}
  \item Submit arbitrary APL source to the pipeline.
  \item Tamper with intermediate artifacts \emph{on disk} between
        boundaries, provided they can forge or delete receipts.
  \item Observe the public Bitcoin ledger and Bifrost chain.
\end{itemize}
We assume the adversary \emph{cannot}:
\begin{itemize}
  \item Break Ed25519 (Plasma Gate) or SHA-256.
  \item Rewrite sufficiently deep Bitcoin history (proof-of-work cost).
  \item Execute arbitrary code inside the Lean kernel's trusted core.
\end{itemize}

\section{Boundary-by-Boundary Resistance}
\label{sec:threat-bounds}

\begin{table}[ht]
\centering
\caption{Adversarial resistance per boundary.}\label{tab:threat}
\small
\begin{tabular}{p{2.6cm} p{5.0cm} p{4.0cm}}
\toprule
\textbf{Boundary} & \textbf{Defends against} & \textbf{Residual} \\
\midrule
Layer 1 $\to$ 2 & malformed APL & none (fails loudly) \\
Layer 3 (Liquid) & out-of-contract numeric values & SMT bound tightness \\
Layer 4 (Lean) & incorrect closed form & trusted core size \\
Layer 5 (C$--$) & unsound translation & compiler correctness \\
Layer 9 (WORM) & silent artifact swap & hash + signature forge \\
Layer 10 & incompleteness drift & human supply of proofs \\
\bottomrule
\end{tabular}
\end{table}

\section{The Forge Game}
\label{sec:threat-forge}

The strongest attack is to forge a receipt that claims a boundary held when
it did not. This requires either forging an Ed25519 signature (assumed
hard) or rewriting the Bifrost block and its Bitcoin anchor (assumed
economically infeasible beyond a small depth). Thus a verifier who checks
the anchor (\cref{ch:anchors}) is protected against even a fully
compromised build machine, provided the anchor transaction is confirmed.

\section{Out-of-Scope Assumptions}
\label{sec:threat-out}

We do \emph{not} claim protection against:
\begin{itemize}
  \item Side channels in the host CPU during C$--$ execution.
  \item Compromise of the Lean trusted core itself (a known,
        unavoidable base).
  \item A broken or colluding Bitcoin majority (the standard
        ledger-security assumption).
\end{itemize}
These are standard limitations of any verified system that ultimately
rests on a small trusted base and an external anchor.

\section{Defense-in-Depth Summary}
\label{sec:threat-summary}

The layered design means an adversary must break \emph{every} relevant
boundary to subvert a result undetectably; a single broken boundary is
caught by the build gate (\cref{ch:build}) and recorded as a failed
receipt rather than a silent pass. This is the practical realization of
the formal model in \cref{ch:formal}.
